\documentclass[11pt,a4paper]{article}

\usepackage{amsmath}
\usepackage{amssymb}
\usepackage{geometry}
\usepackage{booktabs}
\geometry{margin=2.5cm}

\title{Monte Carlo Pricing of Asian Options}
\author{Navy}
\date{\today}

\begin{document}
\maketitle

\section{Setup}

Take $S_t$ to follow geometric Brownian motion under the risk-neutral measure
$\mathbb{Q}$:
%
\begin{equation}
    dS_t = r S_t \, dt + \sigma S_t \, dW_t ,
\end{equation}
%
with $r$ the risk-free rate, $\sigma$ the volatility and $W_t$ a standard
Brownian motion. The solution is
%
\begin{equation}
    S_T = S_0 \exp\!\left( \left(r - \tfrac{1}{2}\sigma^2\right) T
          + \sigma \sqrt{T}\, Z \right),
    \qquad Z \sim \mathcal{N}(0,1),
    \label{eq:gbm}
\end{equation}
%
and equation~\eqref{eq:gbm} is what the simulation samples from directly.

\section{European call as a control}

The European call has a closed form,
%
\begin{equation}
    C = S_0 \Phi(d_1) - K e^{-rT} \Phi(d_2),
    \qquad
    d_{1,2} = \frac{\ln(S_0/K) + \left(r \pm \tfrac{1}{2}\sigma^2\right)T}
                   {\sigma\sqrt{T}},
\end{equation}
%
where $\Phi$ is the standard normal CDF. Since the answer is already known,
pricing it by simulation is a check on the estimator. If the simulated price sits inside the confidence interval
around the analytic price, the sampling and discounting are behaving.

\section{Asian call}

An arithmetic-average Asian call pays
%
\begin{equation}
    \max\!\left( \frac{1}{n}\sum_{i=1}^{n} S_{t_i} - K,\; 0 \right).
\end{equation}
%
A sum of lognormal variables is not lognormal, so there is no closed form to
fall back on. The estimator is
%
\begin{equation}
    C_{\text{Asian}} = e^{-rT}\, \mathbb{E}^{\mathbb{Q}}
        \!\left[ \max\!\left( \bar{S} - K,\, 0 \right) \right]
    \approx \frac{e^{-rT}}{M} \sum_{j=1}^{M}
        \max\!\left( \bar{S}^{(j)} - K,\, 0 \right)
\end{equation}
%
over $M$ paths. Convergence is $O(M^{-1/2})$, so accuracy
is bought expensively, and cutting the error in half costs four times the paths.

\section{Results}

With $S_0 = K = 100$, $r = 5\%$, $\sigma = 20\%$ and $T = 1$:

\begin{center}
\begin{tabular}{lrr}
\toprule
 & \textbf{Price} & \textbf{95\% interval} \\
\midrule
European, analytic     & 10.4506 & n/a \\
European, Monte Carlo  & 10.4520 & $\pm\,0.0204$ \\
Asian, Monte Carlo     &  5.7661 & $\pm\,0.0350$ \\
\bottomrule
\end{tabular}
\end{center}

The simulated European price differs from the analytic one by $0.0014$, well
inside the interval, so the estimator is sound. The Asian call prices at
roughly half the European. Averaging over the path damps the variance of the
terminal quantity, and a less volatile underlying quantity carries less
optionality.

\end{document}
